\GalSourcePageBoundary{087}{L0086}{58}{58}
$i$~étant une racine primitive de
\[
x^{p^{2}-1}\! - 1 = 0\ (\mod p),
\]
il est clair que toute substitution de la forme
\[
[a_{k_{1}+k_{2} i},\ a_{(m_{1}+m_{2}i)(k_{1}+k_{2}i)}]
\]
sera une substitution linéaire; \GalSourceError{mais, dans ces substitutions, aucune
lettre ne reste à la même place, et elles sont au nombre
de~$p^{2} - 1$.}{W10-SE002}

Nous avons donc un système de $p^{2} - 1$ permutations tel que,
dans chacune de ses substitutions, toutes les lettres varient, à
l'exception de~$a_{0,0}$. Combinant ces substitutions avec les $p^{2} - p$
dont il est parlé plus haut, nous aurons
\[
(p^{2} - 1)(p^{2} - p) \text{ substitutions}.
\]

Or, nous avons vu \textit{a~priori} que le nombre des substitutions
où $a_{0,0}$~reste fixe ne pouvait être plus grand que $(p^{2} - 1)(p^{2} - p)$.
Donc il est précisément égal à $(p^{2} - 1)(p^{2} - p)$, et le groupe
linéaire total aura en tout
\[
p^{2}(p^{2} - 1)(p^{2} - p) \text{ permutations}.
\]

Il reste à chercher les diviseurs de ce groupe, qui peuvent
jouir de la propriété d'être solubles par radicaux. Pour cela, nous
allons faire une transformation qui a pour but d'abaisser autant
que possible les équations générales de degré~$p^{2}$ dont le groupe
serait linéaire.

Premièrement, comme les substitutions circulaires d'un pareil
groupe sont telles, que toute autre substitution du groupe les
transforme les unes dans les autres, on pourra abaisser l'équation
d'un degré et considérer une équation de degré $p^{2} - 1$ dont le
groupe n'aurait que des substitutions de la forme
\[
(b_{k_{1}, k_{2}},\ b_{m_{1}k_{1}+n_{1}k_{2},\: m_{2}k_{1}+n_{2}k_{2}}),
\]
les $p^{2} - 1$ lettres étant
\[
\begin{array}{*{5}{l}}
         &b_{0,1}, &b_{0,2}, &b_{0,3}, &\dots,\\
b_{1,0}, &b_{1,1}, &b_{1,2}, &b_{1,3}, &\dots,\\
b_{2,0}, &b_{2,1}, &b_{2,2}, &b_{2,3}, &\dots,\\
\dots,  &\dots,  &\dots,  &\dots,  &\dots.
\end{array}
\]
